\section{Finite-export residual accounting}
\label{sec:orf-export-errors}
An admissible factor need not give exactly proportional binary64 products.
The following elementary bound makes the required reserve explicit. The selector
extension uses it during selection; earlier experimental policies are unchanged. Let
$\widetilde a=\alpha a+e_a$ and $\widetilde b=\alpha b+e_b$ be the actual
written row, with $\alpha>0$. If its observed violation is at most $\epsilon$,
then
\begin{equation}\label{eq:orf-export-transfer}
 (a^Tx-b)_+\le
 \frac{\epsilon+|e_a^Tx|+|e_b|}{\alpha}.
\end{equation}
Indeed, $\alpha(a^Tx-b)=\widetilde a^Tx-\widetilde b-e_a^Tx+e_b$;
taking positive parts and using the triangle inequality proves the bound.
The same argument with absolute values covers equality rows.

If known finite bounds give $|x_j|\le X_j$ wherever $e_{a,j}\ne0$, set
$B(\alpha)=\sum_j|e_{a,j}|X_j+|e_b|$. With the original-unit integer-rounding
allowance $E=\sum_{j\in J}|a_j|\eta_j$, a sufficient condition is
\begin{equation}\label{eq:orf-export-contract}
 \epsilon+B(\alpha)+\alpha E\le\alpha\delta.
\end{equation}
This accounts for the products and decimal serialization actually imported by
the solver, rather than only for a nominal floating-point multiplication.
Errors can be computed as exact rational differences of the input and written
binary64 data. Because $B$ depends on the selected factor, this is a candidate
check, not a claim that arbitrary-factor feasibility is the same closed
interval as in the main text's real-factor characterization. Exact dyadic export has $B=0$ and recovers
the previous contract. Without suitable finite bounds, the pointwise form
\eqref{eq:orf-export-transfer} remains available after solving; no uniform
reserve is asserted. Neither form validates a dual bound or replaces the
original-point audit.
